% 原创教学几何示意；三幅图均使用固定坐标。
\begin{tikzpicture}[x=1mm,y=1mm,font=\footnotesize]
  \foreach \shift/\title in {0/靠近中心,57/留出分类间隔,114/比较相对远近}{
    \begin{scope}[xshift=\shift mm]
      \node[font=\heiti\small] at (22,44) {\title};
      \fill[BookPaper] (0,0) rectangle (46,36);
    \end{scope}
  }
  \foreach \x/\y in {5/6,8/17,16/7,20/14}{
    \draw[FigureInput!45] (\x,\y)--(12,11);
    \fill[FigureInput] (\x,\y) circle (1);
  }
  \foreach \x/\y in {26/22,31/31,40/25,38/18}{
    \draw[FigureLoss!50] (\x,\y)--(34,24);
    \fill[FigureLoss] (\x-0.9,\y-0.9) rectangle (\x+0.9,\y+0.9);
  }
  \foreach \x/\y in {12/11,34/24}{
    \draw[FigureModel,very thick] (\x-2,\y)--(\x+2,\y) (\x,\y-2)--(\x,\y+2);
  }
  \begin{scope}[xshift=57mm]
    \fill[FigureModel!10] (18,0) rectangle (28,36);
    \draw[FigureModel,thick] (23,0)--(23,36);
    \draw[FigureModel,dashed] (18,0)--(18,36) (28,0)--(28,36);
    \foreach \x/\y in {6/8,10/22,15/30,18/15}{
      \fill[FigureInput] (\x,\y) circle (1);
    }
    \foreach \x/\y in {28/9,32/25,37/17,40/31}{
      \fill[FigureLoss] (\x-0.9,\y-0.9) rectangle (\x+0.9,\y+0.9);
    }
    \draw[<->,BookMuted] (18,39)--(28,39);
  \end{scope}
  \begin{scope}[xshift=114mm]
    \draw[FigureModel,thick] (9,10)--(18,24);
    \draw[FigureLoss,thick,dashed] (9,10)--(37,17);
    \fill[FigureInput] (9,10) circle (1.4);
    \fill[FigureModel] (18,24) circle (1.4);
    \fill[FigureLoss] (36,16) rectangle (38,18);
    \node[below] at (9,8) {$\bm a$};
    \node[above] at (18,26) {$\bm p$};
    \node[above] at (37,19) {$\bm n$};
  \end{scope}
  \node at (23,-7) {样本与原型};
  \node at (80,-7) {样本与边界};
  \node at (137,-7) {样本与样本};
\end{tikzpicture}
